\documentclass[11pt,a4paper]{article}

% ============================================================
% PACKAGES
% ============================================================
\usepackage[margin=2.3cm]{geometry}
\usepackage{mathptmx}
\usepackage[T1]{fontenc}
\usepackage{amsmath,amssymb}
\usepackage{booktabs}
\usepackage{array}
\usepackage{tabularx}
\usepackage{microtype}

% ============================================================
% STYLE
% ============================================================
\setlength{\parindent}{0pt}
\setlength{\parskip}{0.55em}
\setlength{\tabcolsep}{5pt}

% ============================================================
% DOCUMENT
% ============================================================
\begin{document}

% ------------------------------------------------------------
% TITLE
% ------------------------------------------------------------
\begin{center}
  {\LARGE\bfseries Vision-Guided Multi-Objective Reinforcement
  Learning\\[0.1em] for UR5 Manipulation}\\[0.4em]
  {\large Project Overview and Expected Results}\\
  \vspace{0.2em}
  {\small Research proposal summary}
\end{center}

\vspace{0.4em}
\noindent\rule{\linewidth}{1pt}
\vspace{0.6em}

% ------------------------------------------------------------
% ABSTRACT
% ------------------------------------------------------------
\begin{abstract}
Real-world robot manipulation requires a robot to achieve its task while
balancing conflicting objectives such as task success, energy consumption,
movement smoothness, and safety. The literature provides each of the
required capabilities --- vision-guided control, multi-objective
reinforcement learning (MORL), 6-DoF manipulation, preference
conditioning, and safety as a constraint --- but no existing system
combines all of them on an articulated arm operating in dynamic scenes.
This project proposes a preference-conditioned, vision-guided
multi-objective policy for a UR5 robot, built on a 3D point-cloud
representation, a multi-objective critic, a diffusion-based visuomotor
policy, and a constrained-MDP formulation of safety. The expected
outcomes, including a target of at least 70\% static success on the real
robot and a hypervolume at or above existing MORL baselines, are
hypotheses to be confirmed by experiment rather than measured results.
\end{abstract}

% ------------------------------------------------------------
% 1. INTRODUCTION
% ------------------------------------------------------------
\section{Introduction and Motivation}

The immediate stimulus for this project is that two otherwise mature
research lines --- multi-objective reinforcement learning and
vision-guided manipulation --- have developed largely independently.
MORL methods provide principled tools for balancing several rewards,
while visuomotor policy learning has delivered strong results on
contact-rich arm control. What is missing is their integration into a
single system with the following five properties:

\begin{enumerate}
  \item \textbf{Vision:} RGB-D / point-cloud observations of the scene.
  \item \textbf{MORL:} vector rewards and a Pareto-optimal solution set.
  \item \textbf{Manipulation:} contact-rich, 6-DoF arm control.
  \item \textbf{Preference conditioning:} runtime modification of
  behavior via a preference vector $\lambda$, without retraining.
  \item \textbf{Safety as a constraint:} collision, velocity, and
  workspace limits expressed as a constrained MDP (CMDP), separate from
  the preference trade-off.
\end{enumerate}

The objective is to obtain one trained policy that observes the scene,
acts on the arm, and whose trade-off among success, energy, and
smoothness can be changed at run time --- all while respecting hard
safety limits.

% ------------------------------------------------------------
% 2. LITERATURE REVIEW
% ------------------------------------------------------------
\section{Literature Review}

The review summarizes the relevant methods in five thematic threads and
then identifies the research gap. Numerical values quoted below are those
reported by the cited publications.

% -------- 2.1 --------
\subsection{Multi-Objective Reinforcement Learning}

The MORL literature supplies the algorithmic basis for policies that
balance multiple objectives.

\textbf{Envelope MOQ} \cite{yang2019} trains a single network
$Q(s,a,\omega)$ over the preference space. Its envelope Bellman operator
and hindsight experience replay provide the foundation for
single-policy preference conditioning.

\textbf{PGMORL} \cite{xu2020} maintains a population of PPO policies
trained under different preference vectors, guided toward the Pareto
front by a Gaussian-process prediction model. Solutions are evaluated by
hypervolume and sparsity.

\textbf{MO-MPO} \cite{abdolmaleki2022} replaces direct scalar weighting
with per-objective KL budgets, yielding a scale-invariant treatment of
multiple objectives.

\textbf{DOL} \cite{mossalam} demonstrates multi-objective learning with
image observations in the ImageDeepSea environment; the visual setting is
a toy grid rather than a camera-based manipulation problem.

\textbf{MO-Gymnasium} \cite{felten2023} provides the standard
vector-reward interface for multi-objective environments, which the
proposed benchmark adopts.

\textit{Limitation.} These methods form the MORL foundation, but none
simultaneously addresses visual 6-DoF manipulation in dynamic scenes with
explicit safety constraints.

% -------- 2.2 --------
\subsection{Vision-Guided Robot Manipulation}

Vision-guided manipulation has advanced rapidly through imitation and
visuomotor policy learning.

\textbf{Diffusion Policy} \cite{chi2023} represents the policy as a
conditional diffusion process over action sequences. Its reported UR5
setup uses $T_o=2$, $T_p=16$, and $T_a=8$, a 10\,Hz policy with 125\,Hz
interpolation, a Cartesian velocity limit of $<0.43$\,m/s, and a
workspace position at least 1\,cm above the table.

\textbf{3D Diffusion Policy (DP3)} \cite{ze2024} learns manipulation from
sparse point clouds of 512 or 1024 points, reporting a 24.2\% relative
gain over image-based baselines across 72 simulation tasks and 85\%
real-robot success with 40 demonstrations.

\textbf{ACT / ALOHA} \cite{zhao2023}, \textbf{Mobile ALOHA} \cite{fu2024},
\textbf{UMI} \cite{chi2024}, and \textbf{IRIS} \cite{jiang2025}
are further demonstration-based visuomotor manipulation systems.

\textit{Limitation.} These approaches primarily optimize a single task
objective and provide no runtime mechanism for changing the trade-off
among multiple objectives.

% -------- 2.3 --------
\subsection{Robot-Scale Multi-Objective RL}

\textbf{GCR-PPO} \cite{munn2025} is the closest algorithmic analogue of
the proposed multi-objective critic. It uses a multi-head critic
$V_{\phi}:S\rightarrow\mathbb{R}^{K}$ with component-wise generalized
advantage estimation, ratio-preserving advantage normalization, and
asymmetric PCGrad. The method reports a $\sim$4.55\% average return gain
over PPO, with gains up to 9.5\% on multi-objective suites.

These values are treated in this project only as an \emph{algorithmic
reference} for the order of magnitude of improvement; they are not adopted
as numerical targets for the UR5 system. GCR-PPO operates on
proprioceptive observations and does not combine a visual encoder,
runtime preference conditioning, or CMDP safety.

% -------- 2.4 --------
\subsection{Safety-Constrained MORL}

\textbf{PRC-CMORL} \cite{he2023} formulates multi-objective robotic
control as a constrained MDP. Its key distinction is that a
\emph{safety objective} is not the same as a \emph{safety constraint}:
if safety enters only as a reward term, the policy may trade safety away
for performance when $\lambda$ changes. A CMDP instead requires the
safety limits to be satisfied explicitly, for example

\begin{equation}
d(s,O)\geq D_{\mathrm{safe}},\qquad
\|\dot{x}\|<0.43~\mathrm{m/s}.
\end{equation}

In the proposed system the preference vector therefore controls only the
task-objective trade-off among success, energy, and smoothness, while
safety is handled separately as a constraint.

% -------- 2.5 --------
\subsection{Closest Intersections: Vision and MORL}

\textbf{Demonstration-enhanced adaptable multi-objective navigation}
\cite{deheuvel2025} is the closest combination of MORL, preference
conditioning, dynamic environments, and sim-to-real transfer. It employs a
preference interpolator $I(\lambda)=\lambda^{p}$ and adapts the
preference at runtime without retraining. However, it uses lidar-based
wheeled navigation rather than visual 6-DoF manipulation, and collision is
treated as a sparse objective rather than a CMDP constraint.

\textbf{POVNav} \cite{pushp2025} combines Pareto reasoning with visual
input and explicit safety constraints, but its focus is classical
navigation planning rather than a learned preference-conditioned
manipulation policy.

% -------- 2.6 --------
\subsection{Research Gap}

The literature thus supplies the individual building blocks described
above, while the combined system remains unaddressed:

\begin{center}
\textbf{Vision + MORL + 6-DoF Manipulation\\
+ Preference Conditioning + CMDP Safety}
\end{center}

on a UR5 arm in static and dynamic environments. The following summary
tables make the gap explicit. Capability comparison:

\renewcommand{\arraystretch}{1.25}
{\small
\begin{tabular}{l c c c c c c}
\toprule
\textbf{Method} & \textbf{Vision} & \textbf{MORL} &
\textbf{Dynamic} & \textbf{Manip.} & \textbf{Preference} &
\textbf{Safety} \\
\midrule
Envelope MOQ (2019) & $\times$ & \checkmark & $\times$ & $\times$
& \checkmark & $\times$ \\
PGMORL (2020) & $\times$ & \checkmark & $\times$ & $\times$
& \checkmark & $\times$ \\
MO-MPO (2022) & $\times$ & \checkmark & $\times$ & partial
& \checkmark & $\times$ \\
Diffusion Policy (2023) & \checkmark & $\times$ & \checkmark &
\checkmark & $\times$ & partial \\
DP3 (2024) & \checkmark & $\times$ & \checkmark & \checkmark &
$\times$ & empirical \\
ACT / ALOHA (2023) & \checkmark & $\times$ & \checkmark &
\checkmark & $\times$ & $\times$ \\
GCR-PPO (2025) & $\times$ & \checkmark & partial & \checkmark &
$\times$ & $\times$ \\
PRC-CMORL (2023) & $\times$ & \checkmark & $\times$ & $\times$ &
\checkmark & \checkmark \\
MORL navigation (2025) & lidar & \checkmark & \checkmark & $\times$ &
\checkmark & $\times$ \\
POVNav (2025) & \checkmark & classical & \checkmark & $\times$ &
partial & \checkmark \\
\textbf{This work} & \textbf{\checkmark} & \textbf{\checkmark} &
\textbf{\checkmark} & \textbf{\checkmark} & \textbf{\checkmark} &
\textbf{\checkmark} \\
\bottomrule
\end{tabular}}

\vspace{0.6em}
What each field provides and what remains:

\renewcommand{\arraystretch}{1.25}
\begin{tabularx}{\linewidth}{@{}l X X@{}}
\toprule
\textbf{Research area} & \textbf{What it provides} &
\textbf{Remaining gap} \\
\midrule
MORL & Pareto optimization & Limited visual manipulation \\
Vision manipulation & Visual 6-DoF control & Mostly single-objective \\
GCR-PPO & Multi-objective critic & No visual preference-conditioned
UR5 \\
PRC-CMORL & Safety constraints & No visual manipulation \\
MORL navigation (2025) & Preference adaptation in dynamic scenes &
Navigation rather than manipulation \\
POVNav & Visual Pareto planning + safety & Classical navigation
planning \\
\textbf{This work} & \textbf{Combines all five} &
\textbf{UR5 dynamic manipulation} \\
\bottomrule
\end{tabularx}

\vspace{0.6em}
The proposed study therefore investigates this combination on a UR5 in
static and dynamic scenes.

% -------- 2.7 --------
\subsection{Positioning Relative to Vision-Guided Baselines}

The vision-guided methods reviewed above are the \emph{visual-manipulation
foundation} (the adopted backbone) of the proposed system, while the MORL
methods form its \emph{algorithmic foundation}. The proposed approach does
not reproduce their single-objective benchmarks; it evaluates on different
criteria. This yields two distinct classes of claims.

\emph{Where the proposed approach can improve on these works:}
\begin{itemize}
  \item \textbf{Dynamic scenes.} The vision-guided works largely target
  static-scene manipulation. The proposed benchmark adds a moving
  obstacle, a moving target, and sudden appearance, which those methods
  do not handle as a learned behavior.
  \item \textbf{One model, many behaviors.} They return a single fixed
  trade-off. The proposed policy produces a Pareto set of behaviors and
  lets the user select among them at run time.
  \item \textbf{Runtime preference, no retraining.} None of the
  vision-guided methods offers preference conditioning; the proposed
  policy changes its success--energy--smoothness trade-off without
  retraining.
  \item \textbf{Explicit safety.} Their safety handling is empirical or
  partial. The CMDP formulation yields explicit guarantees, evaluated as
  measurable violations (static $\approx 0\%$).
  \item \textbf{Controlled attribution.} The proposed evaluation isolates
  the effect of the multi-objective critic, the preference input, and the
  safety constraint (C3), which the vision-guided works do not provide.
\end{itemize}

\emph{Where the proposed approach does not claim to outperform them:}
\begin{itemize}
  \item \textbf{Raw static success.} DP3 is adopted as the backbone, so
  beating its 85\% real success is not claimed; the static target of
  $\geq 70\%$ is deliberately lower because the system also optimizes
  energy, smoothness, and safety.
  \item \textbf{Representation gain.} The $+24.2\%$ point-cloud advantage
  is inherited from the backbone, not improved.
  \item \textbf{Cross-benchmark comparability.} Different tasks, sensors,
  and data scale prevent a direct numerical comparison with their
  published numbers.
\end{itemize}

In short, the defensible claim is: \emph{better under multi-objective,
preference, and safety criteria on a dynamic benchmark --- not higher
single-objective success than the vision baselines}.

% -------- 2.8 --------
\subsection{Overall Contribution: From MORL to This Work}

The overall contribution of the paper is best explained as a chain of
three stages: the field of multi-objective reinforcement learning, its
robot-scale realization, and the proposed unification. Each stage
contributes a capability and leaves a limitation that the next stage
addresses.

\emph{Stage 1 --- Contribution of MORL.}
The MORL field supplies the principled formulation of multi-objective
control: vector rewards, Pareto-optimal solution sets, and
preference-conditioned policies (Envelope MOQ, PGMORL, MO-MPO), together
with a standard benchmark interface (MO-Gymnasium).
\textit{Limitation carried forward:} these results are demonstrated
mostly in abstract or toy settings, without visual 6-DoF manipulation,
dynamic scenes, or explicit constraints.

\emph{Stage 2 --- Contribution of robot-scale MORL.}
This stage scales MORL to realistic robot-learning problems. GCR-PPO
provides a multi-head critic (per-objective value heads, component-wise
GAE, ratio-preserving advantage normalization, asymmetric PCGrad) with
$\sim$4.55\% average and up to 9.5\% return gains, and the IROS 2025
navigation work demonstrates runtime preference adaptation.
\textit{Limitation carried forward:} proprioceptive or lidar inputs only;
no vision-guided manipulation; collision treated as an objective rather
than a CMDP constraint.

\emph{Stage 3 --- Contribution of this work.}
The paper unifies the two stages into a single system:
\begin{itemize}
  \item \textbf{C1:} a preference-conditioned, vision-guided MORL policy
  --- DP3 point-cloud backbone, multi-objective critic, and diffusion
  policy --- with runtime $\lambda$ and no retraining;
  \item \textbf{C2:} a multi-objective UR5 benchmark (static scene plus
  three dynamic scenes) wrapped in MO-Gymnasium;
  \item \textbf{C3:} a controlled evaluation that separates the effects
  of multi-objective learning, preference conditioning, the visual
  representation, and safety, together with an explicit sim-to-real study.
\end{itemize}

\emph{Limitation resolved:} the UR5 static-and-dynamic manipulation
setting that the earlier stages leave open.

The contribution chain is summarized below.

\renewcommand{\arraystretch}{1.3}
\begin{tabularx}{\linewidth}{@{}>{\raggedright\arraybackslash}X
>{\raggedright\arraybackslash}X >{\raggedright\arraybackslash}X@{}}
\toprule
\textbf{Stage} & \textbf{Contribution} &
\textbf{Limitation carried forward} \\
\midrule
MORL field & Multi-objective formulation, Pareto sets, preference tools
& Not demonstrated on visual robots \\
Robot-scale MORL & Critic/GAE/PCGrad machinery; runtime preference
adaptation & No vision-guided manipulation; safety as objective \\
\textbf{This work} & \textbf{Unify MORL + robot scale + vision + CMDP on
a UR5} & \textbf{UR5 static--dynamic manipulation (the gap)} \\
\bottomrule
\end{tabularx}

% ------------------------------------------------------------
% 3. PROPOSED METHOD
% ------------------------------------------------------------
\section{Proposed Method}

\begin{enumerate}
  \item \textbf{Visual representation.} A DP3-style point-cloud encoder
  reduces the RGB-D observation to 512 points. The policy is warm-started
  from 10--20 demonstrations per task.
  \item \textbf{Multi-objective critic.} One value head per objective,
  trained with per-objective generalized advantage estimation,
  ratio-preserving advantage normalization, and asymmetric PCGrad to
  mitigate conflicting gradients.
  \item \textbf{Preference-conditioned policy.} A diffusion-based
  visuomotor policy conditioned on the runtime preference vector
  $\lambda=(\lambda_{\mathrm{success}},\lambda_{\mathrm{energy}},
  \lambda_{\mathrm{smoothness}})$, so the trade-off can be changed
  without retraining.
  \item \textbf{CMDP safety.} Constraints $d(s,O)\geq D_{\mathrm{safe}}$
  and $\|\dot{x}\|<0.43$\,m/s are enforced separately from
  $\lambda$.
  \item \textbf{Benchmark.} A new UR5 multi-objective benchmark following
  the MO-Gymnasium interface, with one static scene and three dynamic
  scenes (moving obstacle, moving target, sudden appearance), so that all
  baselines share the same interface and conditions.
  \item \textbf{Evaluation.} Simulation and real-robot training are run
  in parallel so that the sim-to-real gap is measured explicitly for every
  reported metric.
\end{enumerate}

% ------------------------------------------------------------
% 4. EXPECTED RESULTS
% ------------------------------------------------------------
\section{Expected Results vs. Reported Results}

The following table places the measured results reported in the
literature beside the expected outcomes of this study.

\renewcommand{\arraystretch}{1.3}
\begin{tabularx}{\linewidth}{@{}>{\raggedright\arraybackslash}X
>{\raggedright\arraybackslash}X >{\raggedright\arraybackslash}X@{}}
\toprule
\textbf{Quantity} & \textbf{Reported result (measured, source)} &
\textbf{Proposed target (expected)} \\
\midrule
Point cloud & 512 or 1024 points (DP3 \cite{ze2024}) & 512 points \\
Demonstrations & 10--20 per task (DP3 \cite{ze2024}) & 10--20
(warm-start) \\
Real-robot success & 85\% with 40 demos (DP3 \cite{ze2024}) &
$\geq 70\%$ (static target) \\
Point cloud vs.\ images & $+24.2\%$ relative (DP3 \cite{ze2024}) &
Representation choice, not a gain target \\
Control loop & 125\,Hz (Diffusion Policy \cite{chi2023}) &
125\,Hz \\
Cartesian velocity & $<0.43$\,m/s (Diffusion Policy \cite{chi2023}) &
$<0.43$\,m/s (constraint) \\
Multi-objective gain & $\sim 4.55\%$ avg / $\le 9.5\%$ max
(GCR-PPO \cite{munn2025}) & Reference only; we report our own values \\
Pareto quality & PGMORL / Envelope MOQ hypervolume
\cite{xu2020,yang2019} & Hypervolume reaches or exceeds those
baselines \\
Safety & CMDP constraint (PRC-CMORL \cite{he2023}) & Static violations
$\approx 0\%$; dynamic lower than baseline \\
Preference change & Retraining or separate policies (baselines) &
3--5 settings, same policy, no retraining \\
Simulation success & --- & $\sim 92\%$ (simulation reference) \\
\bottomrule
\end{tabularx}

\vspace{0.6em}
\textbf{Note.} Values in the ``Proposed target'' column are expected
outcomes of this study and must be confirmed by experiment; they are not
measured results. After the experimental phase, each expected value will
be replaced by its measured counterpart.

% ------------------------------------------------------------
% 5. RESEARCH QUESTIONS
% ------------------------------------------------------------
\section{Research Questions and Expected Outcomes}

\renewcommand{\arraystretch}{1.3}
\begin{tabularx}{\linewidth}{@{}>{\raggedright\arraybackslash}X
>{\raggedright\arraybackslash}X >{\raggedright\arraybackslash}X@{}}
\toprule
\textbf{Question} & \textbf{Comparison} & \textbf{Expected outcome} \\
\midrule
Q1: Can the visual representation support learning?
& Proposed system vs.\ single-objective DP3 / PPO &
$\sim 92\%$ in simulation; $\geq 70\%$ real static success \\
Q2: Does the multi-objective critic help?
& Multi-objective critic vs.\ scalarized (single-number) PPO &
Higher hypervolume / better Pareto front; not a fixed percentage gain \\
Q3: Does runtime preference conditioning work?
& Continuous $\lambda$ vs.\ fixed $\lambda$ &
Same policy, measurably different behavior, no retraining \\
Q4: Does CMDP safety reduce violations?
& CMDP vs.\ collision-as-reward &
Static violations $\approx 0\%$; dynamic violations strictly below the
reward baseline \\
\bottomrule
\end{tabularx}

% ------------------------------------------------------------
% 6. EXPERIMENTAL PLAN
% ------------------------------------------------------------
\section{Experimental Plan (UR5)}

\begin{enumerate}
  \item Control cadence of 125\,Hz with the policy updated at 10\,Hz and
  interpolated; Cartesian velocity capped below 0.43\,m/s.
  \item Three to five preference settings exercised on the real arm,
  without retraining between settings.
  \item Real static success target of at least 70\%.
  \item Success, energy, smoothness, and safety violations reported
  separately for the real robot.
  \item Sim-to-real differences reported per metric (success, energy,
  smoothness, safety).
\end{enumerate}

% ------------------------------------------------------------
% 7. EVALUATION CRITERIA
% ------------------------------------------------------------
\section{Evaluation Criteria}

\begin{itemize}
  \item \textbf{Visual learning.} $\sim 92\%$ simulation and
  $\geq 70\%$ real static success.
  \item \textbf{MORL benefit.} Hypervolume above the scalarized baseline.
  \item \textbf{Runtime preference.} 3--5 $\lambda$ settings produce
  measurably different behavior from a single trained policy.
  \item \textbf{Pareto quality.} Hypervolume at or above PGMORL and
  Envelope MOQ on the same benchmark.
  \item \textbf{Safety.} Static violations $\approx 0\%$; dynamic
  violations strictly below the collision-as-objective baseline.
\end{itemize}

% ------------------------------------------------------------
% 8. REFERENCES
% ------------------------------------------------------------
\section{References}

\begin{thebibliography}{15}

\bibitem{yang2019}
R.~Yang \emph{et al.},
``A Generalized Algorithm for Multi-Objective Reinforcement Learning
and Policy Adaptation,''
NeurIPS, 2019.

\bibitem{xu2020}
J.~Xu \emph{et al.},
``Prediction-Guided Multi-Objective Reinforcement Learning for
Continuous Robot Control,''
ICML, 2020.

\bibitem{abdolmaleki2022}
A.~Abdolmaleki \emph{et al.},
``A Distributional View on Multi-Objective Policy Optimization,''
NeurIPS, 2022.

\bibitem{mossalam}
H.~Mossalam \emph{et al.},
``Multi-Objective Deep Reinforcement Learning.''

\bibitem{felten2023}
F.~Felten \emph{et al.},
``MO-Gymnasium: A Toolkit for Reliable Benchmarking and Research in
Multi-Objective Reinforcement Learning,''
NeurIPS, 2023.

\bibitem{chi2023}
C.~Chi \emph{et al.},
``Diffusion Policy: Visuomotor Policy Learning via Action Diffusion,''
RSS, 2023; IJRR, 2024.

\bibitem{ze2024}
Y.~Ze \emph{et al.},
``3D Diffusion Policy: Generalizable Visuomotor Policy Learning via
Simple 3D Representations,''
RSS, 2024.

\bibitem{zhao2023}
T.~Zhao \emph{et al.},
``Learning Fine-Grained Bimanual Manipulation with Low-Cost Hardware
(ALOHA / ACT),''
RSS, 2023.

\bibitem{fu2024}
Z.~Fu \emph{et al.},
``Mobile ALOHA: Learning Bimanual Mobile Manipulation with Low-Cost
Whole-Body Teleoperation,''
CoRL, 2024.

\bibitem{chi2024}
C.~Chi \emph{et al.},
``Universal Manipulation Interface (UMI),''
RSS, 2024.

\bibitem{jiang2025}
Y.~Jiang \emph{et al.},
``IRIS: An Immersive Robot Interaction System,''
CoRL, 2025.

\bibitem{munn2025}
M.~Munn \emph{et al.},
``Scalable Multi-Objective Robot RL through Gradient Conflict
Resolution,''
arXiv:2509.14816, 2025.

\bibitem{he2023}
X.~He \emph{et al.},
``Personalized Robotic Control via Constrained Multi-Objective
Reinforcement Learning,''
Neurocomputing, 2023.

\bibitem{deheuvel2025}
J.~de Heuvel \emph{et al.},
``Demonstration-Enhanced Adaptable Multi-Objective Robot Navigation,''
IROS, 2025; arXiv:2404.04857.

\bibitem{pushp2025}
D.~Pushp \emph{et al.},
``Navigating the Wild: Pareto-Optimal Visual Decision-Making in Image
Space,''
arXiv:2511.07750, 2025.

\end{thebibliography}

\end{document}